\documentclass[11pt]{article}

\usepackage[T1]{fontenc}
\usepackage[utf8]{inputenc}
\usepackage{lmodern}
\usepackage{microtype}
\usepackage[a4paper,margin=1.15in]{geometry}
\usepackage{amsmath,amssymb,amsthm}
\usepackage{url}
\urlstyle{tt}
\usepackage{xcolor}
\usepackage{hyperref}
\hypersetup{
  colorlinks=true,
  linkcolor=black,
  citecolor=black,
  urlcolor=[rgb]{0.10,0.25,0.55},
  pdftitle={A cut and independent-set bound for sparse halves},
  pdfauthor={Pablo Miguel Argudo},
  pdfsubject={Erdos problem 128: sparse halves in triangle-free graphs},
  pdfkeywords={extremal graph theory, triangle-free graphs, sparse halves, Erdos 128}
}

% One shared counter, numbered by section, as in the papers cited.
\newtheorem{theorem}{Theorem}[section]
\newtheorem{proposition}[theorem]{Proposition}
\newtheorem{lemma}[theorem]{Lemma}
\theoremstyle{remark}
\newtheorem{remark}[theorem]{Remark}
\newcommand{\bh}{\beta_*}
\title{A cut and independent-set bound for sparse halves}
\author{Pablo Miguel Argudo}
\date{9 September 2026}
\begin{document}
\maketitle
\begin{abstract}
We derive the partial bound $2587/100000=0.02587$ for sparse halves of every
finite triangle-free graph, assuming the Balogh--Clemen--Lidick\'y max-cut
theorem. The argument is elementary: no flag algebra or semidefinite
certificate is used. Its one new ingredient is a combination---apply the
max-cut theorem, delete the edges inside the smaller side so that it becomes
independent, and feed the result to an independent-set completion bound. The
completion bound itself is Razborov's Section~4.5 with a single step changed;
Section~2 states precisely what is his and what is not. All scalar constants
are rational and exactly certified. This is not a solution of Erd\H{o}s
problem~128, and no claim of priority is made.
\end{abstract}

\section{Statement and normalization}
For $n>0$ let
\[
 \bh(G)=\min\Big\{n^{-2}\!\!\sum_{uv\in E(G)}\!\!z_uz_v \;:\;
   0\le z_v\le 1,\ \sum_v z_v=n/2\Big\}.
\]

\begin{proposition}[rounding]\label{round}
Every graph $G$ on $n>0$ vertices has a set of exactly $\lfloor n/2\rfloor$
vertices spanning at most $\bh(G)n^2$ edges.
\end{proposition}
\begin{proof}
Fix a minimizer and two coordinates $z_u,z_v\in(0,1)$. Transferring $s$ from
one to the other changes the objective by an affine function of $s$ when
$uv\notin E$, and by one with $s^2$-coefficient $-1$ when $uv\in E$; in both
cases the cost is concave, so some endpoint of the feasible range of $s$ is at
least as good, and there $z_u$ or $z_v$ leaves $(0,1)$. Iterating gives a
minimizer with at most one coordinate in $(0,1)$. If $n$ is even that
coordinate must be $0$ and the minimizer is integral. If $n$ is odd it equals
$1/2$; deleting it removes only non-negative terms and leaves exactly
$\lfloor n/2\rfloor$ selected vertices.
\end{proof}

\begin{proposition}[blow-ups]\label{blow}
Let $G[t]$ be the balanced $t$-fold blow-up. Then $\bh(G[t])=\bh(G)$ and
$D_2(G[t])=t^2D_2(G)$, where $D_2$ is the minimum number of edges whose
deletion makes the graph bipartite.
\end{proposition}
\begin{proof}
Each twin class is independent, so the objective on $G[t]$ depends on a
profile only through the class sums $t w_u=\sum_{i\in u}z_i$; averaging inside
classes leaves it unchanged and gives a profile on $G$ of the same normalized
cost, and lifting is the converse. For $D_2$, a bipartition of $G[t]$ splitting
class $u$ into $a_u$ and $t-a_u$ has internal cost
$\sum_{uv\in E(G)}\big(a_ua_v+(t-a_u)(t-a_v)\big)$, which is affine in each
$a_u$ separately; hence a minimum is attained with every $a_u\in\{0,t\}$, i.e.\
at a lifted bipartition of $G$, of cost $t^2D_2(G)$.
\end{proof}

\begin{theorem}\label{main}
Assume the Balogh--Clemen--Lidick\'y theorem \cite[Theorem~2(a)]{bcl}: for
$n$ large enough, every triangle-free graph satisfies $D_2(G)\le n^2/23.5$.
Then every finite triangle-free graph $G$ has a set of $\lfloor n/2\rfloor$
vertices spanning at most $2587n^2/100000$ edges, strictly so for $n>0$.
\end{theorem}
\noindent For $n=0$ the empty set gives the non-strict statement. The target
$n^2/50$ remains unproved; so does any bound below $2587/100000$.

\begin{remark}
By Proposition \ref{blow} the hypothesis ``$n$ large enough'' in \cite{bcl}
costs nothing: given any triangle-free $G$ of order $n$, apply the cited bound
to $G[t]$ for $t$ large and divide by $t^2$ to get $D_2(G)\le 2n^2/47$. We use
$I=2/47$ throughout. Similarly, Proposition \ref{round} lets us work with the
relaxation $\bh$ and only round at the very end.
\end{remark}

\section{The completion envelope, and what is new in it}
For $\alpha\in[1/3,1/2]$ put
\[
 g_\alpha(p)=2\big(\tfrac12-\alpha\big)(p-\alpha)+p\big(\tfrac12-p\big),
 \qquad
 f(\alpha)=\tfrac12\max_{p\in[\alpha,1/2]}g_\alpha(p),
\]
and $f(\alpha)=0$ for $\alpha\ge1/2$. Since $g_\alpha$ is concave with
unconstrained maximum at $p=3/4-\alpha$, and $3/4-\alpha\le\alpha$ exactly when
$\alpha\ge3/8$,
\begin{equation}\label{fdef}
 f(\alpha)=\begin{cases}
  \frac32\alpha^2-\frac54\alpha+\frac9{32}, & \tfrac13\le\alpha\le\tfrac38
    \quad\text{(interior maximum at }p=\tfrac34-\alpha),\\[2pt]
  \frac{\alpha}{2}\big(\frac12-\alpha\big), & \tfrac38\le\alpha\le\tfrac12
    \quad\text{(maximum at the endpoint }p=\alpha).
 \end{cases}
\end{equation}
$f$ is continuous at $3/8$, where both branches equal $3/128$, and decreasing
on $[1/3,1/2]$, with $f(1/3)=1/32$ and $f(2/5)=1/50$. Also
$f(\alpha)-\frac{\alpha}{2}(\frac12-\alpha)=2(\alpha-\frac38)^2\ge0$, so
$\frac{\alpha}{2}(\frac12-\alpha)\le f(\alpha)$ on the whole interval.

\begin{lemma}\label{ind}
Let $G$ be triangle-free with normalized independence number
$\alpha=\alpha(G)$. If $\alpha\ge1/3$ then $\bh(G)\le f(\alpha)$. Consequently
$\bh(G)\le f(x)$ whenever $G$ has an independent set of size at least $xn$
with $1/3\le x\le1/2$.
\end{lemma}

\noindent\textbf{Attribution.} For $\alpha\ge3/8$ this is exactly Razborov
\cite[Theorem~3.6]{raz}, whose proof occupies his Section~4.5. Our statement
differs from his only in the range $\alpha\in[1/3,3/8)$, and the proof differs
from his in exactly one step:
\begin{itemize}
\item \emph{Case 1 (the changed step).} Razborov derives the concave bound
$2\bh\le g_\alpha(p_b)$ and observes that its maximum sits at
$p_b=3/4-\alpha$, ``which is $\le\alpha$ since we assumed $\alpha\ge3/8$'';
he therefore substitutes $p_b:=\alpha$. For $\alpha<3/8$ that substitution is
unavailable, but the maximum is then interior to $[\alpha,1/2]$ and one may
simply evaluate $g_\alpha$ there. This is the whole of the extension, and it
is a one-line observation on his own display.
\item \emph{Case 2 (unchanged).} His Case 2 uses $\alpha\ge3/8$ only in its
final step, where the requirement is
$\partial Q_1/\partial p_b|_{p_b=\alpha}=(1-3\alpha)/2\le0$, i.e.\
$\alpha\ge1/3$. It therefore applies verbatim on the wider range. We reproduce
it below for self-containedness, not for novelty.
\end{itemize}
The genuinely new content of this note is not Lemma \ref{ind} but its
combination with a max-cut theorem in Section~3.

\begin{proof}[Proof of Lemma \ref{ind}]
The last sentence follows from the first because $\alpha\ge x$ and $f$ is
decreasing. If $\alpha\ge1/2$ then $\bh(G)=0$. By Proposition \ref{blow} we
may assume $n$ even. Densities: all vertex counts and degrees are divided by
$n$, all edge counts by $n^2$, and each edge is counted once.

Let $A$ be a maximum independent set, $k=\frac12-\alpha$, and grow $B\supseteq
A$ by repeatedly adding a vertex with at most $k$ neighbours in the current
set, stopping at size $\frac12$ or when no such vertex exists. In the first
case the cost is at most $k^2\le\alpha k/2\le f(\alpha)$, using $\alpha\ge1/3$
for the middle inequality. Otherwise we obtain $B$ with
\begin{equation}\label{Bprops}
 b\in[\alpha,\tfrac12),\qquad e(B)\le k(b-\alpha),\qquad
 e_B(v)>k\ \ \text{for every } v\notin B .
\end{equation}
Vertices outside $B$ with more than $b/2$ neighbours in $B$ are pairwise
non-adjacent---two of them share a neighbour in $B$---so they form an
independent set $C$ with $|C|/n\le\alpha$. Hence we may pick $D$ disjoint from
$B\cup C$ with $d=1-\alpha-b$, and every $z\in D$ has $k<s_z\le b/2$, where
$s_z$ is its normalized degree into $B$. Write $h=\frac12-b>0$ and note
$d-h=k>0$.

\emph{Case 1: some $v\in B$ has at least $hn$ neighbours in $D$.} Pick $hn$ of
them; they lie in $N(v)$, hence are independent, and lie in $D$, hence send at
most $b/2$ each into $B$. Their union with $B$ has size $b+h=\frac12$ and cost
at most
\[
 H=k(b-\alpha)+\tfrac b2 h=\tfrac12 g_\alpha(b)\le f(\alpha),
\]
the last step by the definition of $f$ as a maximum over $p\in[\alpha,1/2]$,
which contains $b$.

\emph{Case 2: every $v\in B$ has fewer than $hn$ neighbours in $D$.} Fix
$v\in B$, set $E=N(v)\cap D$ and $r=|E|/n<h$, and take the profile equal to
$1$ on $B\cup E$, to $w=(h-r)/(d-r)\in[0,1]$ on $D\setminus E$, and to $0$
elsewhere; it sums to $b+r+(h-r)=\frac12$. Accounting for the six kinds of
pair---inside $B$, $B$ to $E$, inside $E$ (empty, as $E\subseteq N(v)$ is
independent), $B$ to $D\setminus E$, $E$ to $D\setminus E$, and inside
$D\setminus E$ (Mantel)---the cost is at most
\[
 k(b-\alpha)+m(B,E)+\tfrac b2(h-r)+r(h-r)+\tfrac{(h-r)^2}{4}.
\]
For Mantel: in a triangle-free graph of order $N>0$ with $m$ edges every edge
has endpoint degree sum at most $N$, so
$4m^2/N\le\sum_v\deg(v)^2\le Nm$ and $m\le N^2/4$.

Now average over $v\in B$. The displayed tail is concave in $r$ (its
$r^2$-coefficient is $-3/4$), so it may be evaluated at $\bar r$. A vertex
$z\in D$ enters $E$ with probability $s_z/b$, whence
$\bar r=\frac1{bn}\sum_{z\in D}s_z\in[0,d/2]$ and, using
$s_z^2\le(k+\frac b2)s_z-\frac{kb}{2}$---which is exactly
$(s_z-k)(s_z-\frac b2)\le0$---
\[
 \mathbb E\,m(B,E)=\tfrac1{bn}\textstyle\sum_{z\in D}s_z^2
   \le\big(k+\tfrac b2\big)\bar r-\tfrac{kd}{2}.
\]
Collecting terms, some half costs at most
\[
 Q=k\big(b-\alpha-\tfrac d2\big)+\tfrac{bh}{2}+\tfrac{h^2}{4}
   +\big(k+\tfrac h2\big)\bar r-\tfrac34\bar r^{\,2}.
\]
$Q$ is concave in $\bar r$ with $\partial Q/\partial\bar r=(b-\alpha)/4\ge0$ at
$\bar r=d/2$, so it increases on $[0,d/2]$ and
\[
 Q\le\tfrac{13}{16}\alpha^2-\tfrac98\alpha b-\tfrac3{16}b^2-\tfrac14\alpha
     +\tfrac12 b
   =\tfrac{\alpha k}{2}+\tfrac{1-3\alpha}{2}(b-\alpha)
     -\tfrac3{16}(b-\alpha)^2 .
\]
Both correction terms are $\le0$ because $\alpha\ge1/3$ and $b\ge\alpha$;
hence $Q\le\alpha k/2\le f(\alpha)$.
\end{proof}

\section{Deleting one side of a cut}
This section is the new step. Let $I=2/47$ and take a bipartition $X,Y$ with
$|X|\le|Y|$ realizing $D_2(G)$, with internal edge densities $q=e(X)/n^2$ and
$j=e(Y)/n^2$, so $q+j\le I$. Put
\[
 c=\frac{n}{2|Y|}\in[\tfrac12,1],\qquad x=\frac{|X|}{n}=1-\frac1{2c}.
\]
The profile equal to $c$ on $Y$ and $0$ elsewhere is feasible and gives
\begin{equation}\label{first}
 \bh(G)\le c^2 j .
\end{equation}
Suppose $c\ge3/4$, equivalently $x\ge1/3$. Delete the edges inside $X$: the
resulting graph is still triangle-free and now has $X$ as an independent set of
density $x\in[1/3,1/2]$, so Lemma \ref{ind} applies to it, and restoring the
deleted edges costs at most $q$ because all weights are at most $1$:
\begin{equation}\label{second}
 \bh(G)\le f(x)+q .
\end{equation}
Adding \eqref{first} to $c^2$ times \eqref{second} and using $q+j\le I$,
\begin{equation}\label{combined}
 \bh(G)\le\frac{c^2\big(I+f(x)\big)}{1+c^2},\qquad c\ge\tfrac34 .
\end{equation}
Neither ingredient is new---\cite{bcl} is cited in \cite{raz}---but we have
not found this combination in the literature. Note that \eqref{first} alone is
weak near $c=1$ and \eqref{second} alone is weak near $c=3/4$; the point is
that the two are strong in complementary regimes, and $\eqref{combined}$
interpolates.

\section{Exact scalar certificate}
Put $T=2587/100000$ and $c_0=7797/10000$. The three intervals below cover
$[\frac12,1]$.

For $c\le c_0$, \eqref{first} and $j\le I$ give
$\bh\le Ic^2\le Ic_0^2$, and
$T-Ic_0^2=\frac{1291}{2350000000}>0$.

For $c_0\le c\le\frac45$ we have $x\le3/8$, so the first branch of
\eqref{fdef} applies, and the numerator of $T-c^2(I+f(x))/(1+c^2)$ is
\[
 P(c)=\Big(T-I-\tfrac{17}{32}\Big)c^2+\tfrac{7c}{8}+T-\tfrac38 ,
\]
concave, hence minimized at an endpoint; both are positive:
\[
 P(c_0)=\frac{312483613}{235000000000000},\qquad
 P\big(\tfrac45\big)=\frac{22649}{117500000}.
\]

For $\frac45\le c\le1$ the second branch applies and the numerator is
$Ac^2-Bc+C$ with
\[
 A=T+\tfrac14-I>0,\quad B=\tfrac38,\quad C=T+\tfrac18,\quad
 4AC-B^2=\frac{442569}{2500000000}>0,
\]
so it is positive on all of $\mathbb R$. Proposition \ref{round} converts the
resulting bound on $\bh$ into the statement of Theorem \ref{main}.

\begin{remark}
$T$ is not a free choice: the supremum over $c\in[\frac12,1]$ of the bound
$\min\{Ic^2,\,c^2(I+f(x))/(1+c^2)\}$ equals $0.0258691\ldots$, attained where
the two branches cross near $c_0$. So $2587/100000$ is this argument's own
optimum, rounded up, and no better constant follows from it without a new
ingredient.
\end{remark}

\section{Verification, scope and limitations}
The identities, the rational margins and the finite instances are recomputed by
two independently written verifiers, and the scalar optimization additionally
carries a machine-checked certificate:
\begin{center}
\begin{tabular}{@{}ll@{}}
\texttt{outputs/verify-cut-independent.py} & standard library only\\
\texttt{outputs/verify-independent.py}     & requires \texttt{sympy}, \texttt{networkx}\\
\texttt{formal/verify.py}                  & compiles and replays \texttt{formal/Scalar.lean}
\end{tabular}
\end{center}
\noindent Both verifiers use exact arithmetic only; no floating-point value
enters any acceptance decision. The Lean~4 development certifies the scalar
optimization of Section~4 and the polynomial identities of Section~2,
conditional on scalar hypotheses supplied as arguments.

Not established here: the conjecture $1/50$; any counterexample; a formal proof
of the graph-theoretic content of Sections 2 and 3; and the Balogh--Clemen--%
Lidick\'y theorem, which is assumed as published. The literature audit
underlying the attribution in Section~2 is recorded in \path{LITERATURE.md};
it is not a substitute for review by a specialist.

\section*{Acknowledgement of method}
This work was carried out with substantial assistance from AI coding and
reasoning agents, which produced and checked the algebra, the rational
certificates and the Lean development, and performed the literature audit
recorded in \path{LITERATURE.md}. The author takes responsibility for the
mathematical content. No claim of priority is made and the result has not yet
been read by a specialist referee.

\begin{thebibliography}{9}
\bibitem{bcl} J. Balogh, F. C. Clemen, B. Lidick\'y,
\emph{Max cuts in triangle-free graphs}, arXiv:2103.14179, Theorem~2(a).
\bibitem{raz} A. A. Razborov,
\emph{More about sparse halves in triangle-free graphs},
arXiv:2104.09406, Theorem~3.6 and Section~4.5.
\end{thebibliography}
\end{document}
